% 'nonacm' to remove any running header
\documentclass[nonacm,sigconf]{acmart}

\usepackage{cite}
\usepackage{amsmath,amsfonts}
\usepackage{algorithmic}
\usepackage{graphicx}
\usepackage{textcomp}
\usepackage[svgnames]{xcolor}
\usepackage{booktabs}
\usepackage{tagging}
\usepackage{enumitem}
\usepackage{minted}
\usepackage{needspace} % for orphan optimization

\usepackage{sc25repro}
\pagenumbering{gobble} % no page numbers
\pagestyle{empty} % suppress acmart running header (no \title/\author set)
\makeatletter\def\shorttitle{}\def\@shorttitle{}\makeatother

\setlist[itemize]{leftmargin=*}
\setlist[enumerate]{leftmargin=*}

\newcommand{\profmini}{\textsf{Mini}}
\newcommand{\profdet}{\textsf{Detailed}}
\newcommand{\textmpiw}{\texttt{MPI\_Wait}}

\newmintinline[code]{text}{bgcolor=Gainsboro!40}

% --- environments for verbatim blocks ---
\newenvironment{bash}
{\VerbatimEnvironment\begin{minted}[frame=lines,framesep=2mm,baselinestretch=0.9,bgcolor=WhiteSmoke,fontsize=\footnotesize,obeytabs=true,tabsize=4]{bash}}
{\end{minted}}

\newenvironment{example}
{\VerbatimEnvironment\begin{minted}[bgcolor=GhostWhite,fontsize=\small,obeytabs=true,tabsize=4]{text}}
{\end{minted}}

% Uncomment/comment the following usetag lines
% if you would like to get an explanation or an example.
% Don't forget to comment them for the final submission.
\usetag{explanation}
\usetag{example}

\begin{document}

\twocolumn[%
	{\begin{center}
				\Huge
				Appendix: Artifact Description/Artifact Evaluation
			\end{center}}
]

%%%%%%%%%%%%%%%%%%%%%%%%%%%%%%%%%%%%%%%%%%%%%%%%%%%%%
%  AD Appendix
%%%%%%%%%%%%%%%%%%%%%%%%%%%%%%%%%%%%%%%%%%%%%%%%%%%%%

\appendixAD

\section{Overview of Contributions and Artifacts}

\subsection{Paper's Main Contributions}

\begin{description}
	\item[$C_1$] \textbf{Low-overhead tracing.} Always-on conventional tracing at \textasciitilde0.8--2\% overhead (vs 18--81\% for TAU/dftracer), with Parquet traces 3.6--44$\times$ smaller than other formats.
	\item[$C_2$] \textbf{Faster post-mortem analysis.} ORCA traces enable post-mortem queries multiple orders of magnitude faster than other analytics-oriented tracers.
	\item[$C_3$] \textbf{In-situ analytics.} 
	Flows with operator pushdown to MPI tier execute at sub-2\% overhead while reducing data movement by 98.5--99.999\%. 
	\emph{Not supported by conventional tracers.}
	\item[$C_4$] \textbf{Agentic debugging.} Steerable telemetry lets an LLM agent localize a months-long 4096-rank AMR anomaly in 27 minutes, with 144$\times$ volume reduction when reverting to baseline. \emph{Not supported by conventional tracers.}
\end{description}

\subsection{Computational Artifacts}

\begin{description}
	\item[$A_1$] ORCA-Augmented AMR Code
	\item[$A_2$] Post-Mortem Query Scripts
	\item[$A_3$] LLM Prompt for Agentic Workflow
\end{description}

\begin{center}
	\begin{tabular}{rll}
		\toprule
		Artifact ID & Contributions       & Related                                \\
		            & Supported           & Paper Elements                         \\
		\midrule
		$A_1$       & $C_1$, $C_3$, $C_4$ & Figs.~5, 6, 8 ($C_1$)                  \\
		            &                     & Fig.~9 and Table~1 ($C_3$)             \\
		            &                     & Fig.~10, live system for $A_3$ ($C_4$) \\
		\midrule
		$A_2$       & $C_2$               & Fig.~7                                 \\
		\midrule
		$A_3$       & $C_4$               & Fig.~10                                \\
		\bottomrule
	\end{tabular}
\end{center}

%%%%%%%%%%%%%%%%%%%%%%%%%%%%%%%%%%%%%%%%%%%%%%%%%%%%%%%%
\section{Artifact Identification}
%%%%%%%%%%%%%%%%%%%%%%%%%%%%%%%%%%%%%%%%%%%%%%%%%%%%%%%%

% ----------------------------------------------
% A1 --- ORCA-Augmented AMR Code
% ----------------------------------------------
\newartifact

\artrel

$A_1$ is the instrumented BSP workload --- Phoebus + Parthenon's Sedov Blast Wave 3D AMR code linked against ORCA --- used as the common substrate for three contributions:

\begin{itemize}
	\item $C_1$ (Figs.~5, 6, 8): running $A_1$ under different tracer configurations produces the runtime overhead, trace size, and architectural ablation measurements.
	\item $C_3$ (Fig.~9): running $A_1$ with the in-situ flows from Table~1 produces the pushdown overhead and data-reduction measurements.
	\item $C_4$ (Fig.~10): $A_1$ is the live system that $A_3$'s LLM agent attaches to and steers.
\end{itemize}

\noindent The artifact is provided with two profiles: \profmini{} and \profdet{}.
\begin{itemize}[topsep=0pt]
	\item \profmini{}: triggers short runs, sufficient to verify artifact function.
	\item \profdet{}: requires 262 nodes on the PDL Wolf cluster and triggers longer experiments, necessary to replicate all results.
\end{itemize}

\needspace{5em}
\artexp
\begin{itemize}
	\item \profmini{}: three timesteps of a simple example program execute, generating Parquet traces via ORCA.
	\item \profdet{}: 2000 timesteps of the Blast Wave AMR code, generating Parquet traces via ORCA in two modes---lightweight and detailed. 
	A baseline run is also triggered for overhead comparison.
	\item Trace sizes: ORCA traces take around 10\,B per record, aligning with Fig.~6's measurements.
\end{itemize}

\arttime

\begin{itemize}
	\item Setup: \textasciitilde30 minutes (building the entire project)
	\item Execution: \textasciitilde1 min for \profmini{}, \textasciitilde2--3 hours for \profdet{} (Baseline and ORCA only)
	\item Analysis: $<$10 minutes
\end{itemize}

\artin

\artinpart{Hardware}

\begin{itemize}
	\item \profmini{}: one multi-core Linux node
	\item \profdet{}: 262 nodes on the PDL Wolf cluster
\end{itemize}

\artinpart{Software}\mbox{}\par

\noindent\textbf{Required dependencies, must be present in the environment.}
Versions used by authors are noted in parentheses, 
other environments may necessitate tweaks to build parameters etc.

\begin{itemize}
	\item A Linux environment (Ubuntu 22.04)
	\item An MPI implementation (MVAPICH2 2.3.7)
	\item C/C++ compiler (GCC 11.4)
	\item CMake (3.22)
	\item Rust toolchain (cargo/rustc, stable 1.92)
	\item TAU (2.34.1)
\end{itemize}

\noindent\textbf{Dependencies built by orca-umbrella} \\
(\url{https://github.com/pdlfs/orca-umbrella})

\vspace{0.5em}
\noindent \emph{ORCA-Augmented AMR Code:}
\begin{itemize}[topsep=0pt]
	\item ORCA --- \url{https://github.com/pdlfs/orca}
	\item Phoebus --- \url{https://github.com/lanl/phoebus}
	\item Parthenon --- \url{https://github.com/parthenon-hpc-lab/parthenon}
	\item amr-tools --- \url{https://github.com/pdlfs/amr-tools}
	\item Kokkos 4.0.01 --- \url{https://github.com/kokkos/kokkos}
	\item HDF5 1.12.2 --- \url{https://github.com/HDFGroup/hdf5}
\end{itemize}

\vspace{0.5em}
\noindent\emph{ORCA runtime / data-path:}
\begin{itemize}[topsep=0pt]
	\item Apache Arrow 19.0.0 --- \url{https://github.com/apache/arrow}
	\item Snappy 1.2.2 --- \url{https://github.com/google/snappy}
	\item Mercury --- \url{https://github.com/mercury-hpc/mercury}
	\item libfabric v2.5.1 --- \url{https://github.com/ofiwg/libfabric}
	\item DuckDB v1.2.1 --- \url{https://github.com/duckdb/duckdb}
	\item yaml-cpp 0.9.0 --- \url{https://github.com/jbeder/yaml-cpp}
	\item PSM (forked) --- \url{https://github.com/pdlfs/psm}
\end{itemize}

\vspace{0.5em}
\noindent\emph{Comparison tracers (for complete results replication):}
\begin{itemize}[topsep=0pt]
	\item DFTracer v2.0.2 (forked) --- \url{https://github.com/anku94/dftracer}
	\item Score-P 9.3 (patched, patch automatically applied during build) --- \url{https://www.vi-hps.org/projects/score-p/}
	\item Caliper (patched, patch automatically applied during build) --- \url{https://github.com/LLNL/Caliper}
\end{itemize}

\needspace{4em}
\artinpart{Datasets / Inputs}
\begin{itemize}
	\item AMR decks: \code{blast_wave_3d.<nranks>.pin} (nranks in 512, 1024, 2048, 4096)
	\item \code{wfopts.yml}: tuned config for Wolf cluster, only needed for \profdet{} profile
\end{itemize}

\artinpart{Installation and Deployment}

Preparation commands (indicative, from Ubuntu 22.04, may vary across environments):

\begin{bash}
# Install basic utilities
sudo apt install -y gcc g++ make cmake autoconf \
  automake libtool pkg-config git

# Install rust toolchain via rustup (non-interactive, stable)
curl --proto '=https' --tlsv1.2 -sSf https://sh.rustup.rs \
  | sh -s -- -y --default-toolchain stable
. "$HOME/.cargo/env"
\end{bash}

Now, orca-umbrella can be built, and will also build the rest of the dependencies in-tree.

\begin{bash}
# Install prefix should be a shared dir on the cluster
OR_PREFIX=/tmp/orca-prefix
OR_BUILD=$OR_PREFIX/build
OR_INSTALL=$OR_PREFIX/install
mkdir -p $OR_PREFIX $OR_BUILD $OR_INSTALL

git clone https://github.com/pdlfs/orca-umbrella.git $OR_PREFIX
cd $OR_BUILD
cmake -DCMAKE_INSTALL_PREFIX=$OR_INSTALL $OR_PREFIX/orca-umbrella
make -j<nproc>
\end{bash}

\noindent Once successfully built, all artifacts will be available in \code{$OR_INSTALL}. 
Run the artifact using scripts in \code{$OR_INSTALL/scripts}.

\begin{itemize}
	\item \profmini{}: run \code{run_amr_test.sh}
	\item \profdet{}: Generate a valid hostfile as per MPI configuration, with at least 262 nodes, and run:
			\begin{bash}
OR_HOSTFILE=/path/to/hostfile \
	$OR_INSTALL/scripts/run_orca_detailed.sh
			\end{bash}
\end{itemize}

\artcomp

\begin{itemize}
	\item The relevant MPI codes are executed under exhaustive MPI/Kokkos trace collection modes, as per the configured scales and timesteps.
	\item \profdet{}: Executes a baseline run and ORCA in lightweight and detailed modes at different scales, and reports overhead.
\end{itemize}

\artout

Each run yields timestep-partitioned Parquet traces in the job directory, organized by schema:

\begin{example}
parquet/
|-- kokkos_events/
|   `-- ts=<timestep_range>/
|       `-- ranks=<rank_range>.parquet
|-- mpi_collectives/
|   `-- ...
|-- mpi_messages/
|   `-- ...
`-- orca_events/
\end{example}

\noindent \profdet{}: Wall-clock times are compared with baseline, and measured runtime overhead of different modes is reported.

% ----------------------------------------------
% A2 --- Post-Mortem Query Scripts
% ----------------------------------------------
\newartifact

\artrel

$A_2$ is the post-mortem analytics suite that executes a 
query suite consisting of three queries, against traces captured by running $A_1$.
The three queries in the suite are: outlier waits, outlier collectives, and timestamp range.
It supports $C_2$ by producing the query-latency comparison in Fig.~7.

\artexp

The query suite runs against all available traces. Row counts in executed queries are displayed for verification and time taken by queries is logged as a CSV.

\arttime

\begin{itemize}
	\item Setup: $<$5 minutes
	\item Execution: $<$1 minute (ORCA only)
	\item Analysis: $<$1 minute
\end{itemize}

\artin

\artinpart{Hardware}
One multi-core Linux node.

\artinpart{Software}
ORCA trace queries: Python 3 environment with Polars.
\begin{itemize}
	\item Python (3.12.11)
	\item Polars (1.39.3)
	\item Pandas (2.3.3)
\end{itemize}

\artinpart{Datasets / Inputs}
ORCA traces generated by the AMR codes in $A_1$.

\artinpart{Installation and Deployment}
All Python scripts will be copied to \code{$OR_INSTALL/scripts/orcaquery} during the orca-umbrella build (see $A_1$). 
Polars may be installed by running
\begin{center}
\code{pip install -r requirements.txt}
\end{center}
(or equivalent for the test environment).

\artcomp

The script automatically detects all available runs in the suite, and runs each query for each available run.

\artout

Matching row counts emitted to stdout for verification of successful query execution, and a CSV + a Pandas dataframe logged to stdout with the query latencies.

% ----------------------------------------------
% A3 --- LLM Prompt for Agentic Workflow
% ----------------------------------------------
\newartifact

\artrel

% A3 is the LLM agent specification --- the prompt, tool wiring (ORCA controller CLI + Polars analytics), and operator-in-the-loop protocol --- used to drive the agentic debugging session against a running 4096-rank instance of A1. It supports $C_4$ by producing the localization timeline in Fig.~10.
$A_3$ is a prompt that is supplied to the LLM agent (Claude Opus 4.5 at experiment time)
to explain how to interact with an active ORCA-augmented BSP code, request telemetry, and schedule
in-situ analyses to explain anomalous behavior.
It supports $C_4$ by executing the localization workflow shown in Fig.~10.

\artexp

\begin{enumerate}
	\item (Warmup) LLM is able to interact with the running code, pause, resume, and request and analyze telemetry.
	\item (Exercise) LLM is able to progressively localize the reproduced performance anomaly.
	\item Common case telemetry is two orders of magnitude less than detailed traces.
\end{enumerate}

Caveat: this artifact demonstrates an agentic debugging workflow, not a fixed benchmark. The paper experiment used Claude Opus 4.5, but evaluators may use a comparably capable current coding agent. Exact command trajectories, runtime, and degree of operator intervention may vary, but the intended outcome is the same qualitative result: the agent uses ORCA to localize the anomaly to the same root-cause region. The anomaly itself is specific to our environment.

\arttime

\begin{itemize}
	\item Setup: 5 minutes
	\item Execution: $<$1 hour
	\item Analysis: $<$1 minute
\end{itemize}

\artin

\artinpart{Hardware}
See $A_1$.

\artinpart{Software}
See $A_1$, plus:
\begin{itemize}
	\item Claude Code with a valid Claude subscription or API key --- \url{https://www.anthropic.com/claude-code}
\end{itemize}

\artinpart{Datasets / Inputs}
A prompt \code|agentic-workflow.md| describing ORCA and documenting the experimental workflow.

\artinpart{Installation and Deployment}

Claude Code installation:
\begin{bash}
	curl -sSL https://claude.ai/install.sh | bash
\end{bash}

Deployment:
\begin{enumerate}
	\item \code|run_orca_agentic.sh| triggers a 4096-rank AMR run with configuration parameters designed to disable anomaly mitigations and exacerbate occurrences. Simulation starts in paused mode.
	\item Claude Code must be started on another tab, and provided the IP/port of the controller interface of the ORCA controller. It is then asked to initiate the warmup exercise.
	\item After the warmup exercise is complete, the agent should be asked to start the main exercise. Minimal and guided interventions may be necessary in case it makes suboptimal steering choices.
\end{enumerate}

\artcomp

\begin{itemize}
	\item (Warmup) Agent interacts with ORCA and understands the Parquet output structure and analysis workflows.
	\item (Exercise) Agent uses high-level telemetry to detect the presence of stragglers, and progressively requests additional telemetry by writing custom in-situ flows to identify the root cause (up to \textmpiw{}, but not beyond).
\end{itemize}

\artout

\begin{itemize}
	\item (Warmup) Manipulation effects are visible in simulation stdout.
	\item (Exercise) Agent reports a complete callflow graph of Kokkos regions up to the anomalous \textmpiw{}.
\end{itemize}

%%%%%%%%%%%%%%%%%%%%%%%%%%%%%%%%%%%%%%%%%%%%%%%%%%%%%
%  AE Appendix
%%%%%%%%%%%%%%%%%%%%%%%%%%%%%%%%%%%%%%%%%%%%%%%%%%%%%
\newpage
\appendixAE

\arteval{1}

% Template reference for A1:
% Artifact Setup:
% Provide instructions for installing and compiling libraries and code.
% Offer guidelines on deploying the code to resources.
%
% Artifact Execution:
% Describe the experiment workflow.
% If encapsulated within a workflow description or equivalent (such as a
% makefile or script), clearly outline the primary tasks and their
% interdependencies. Detail the main steps in the workflow. Merely instructing
% to ``Run script.sh'' is inadequate.
%
% Artifact Analysis:
% Provide a description of the expected results and a methodology for
% evaluating these results. Explain how the expected results from the experiment
% workflow correlate with the contributions stated in the article. For example,
% if the article presents results in a figure, the artifact evaluation should
% also produce a similar figure, depicting the same generalizable outcome.
% Authors must focus on these aspects to reduce the time required for others to
% understand and verify an artifact.

\artin

We provide instructions for evaluating ORCA functionality using a single Linux
node. We recommend Ubuntu 22.04 or 24.04 to build and run the artifact, as
minor changes in package names can break compatibility.

First, C++ and Rust toolchains must be available, as well as an MPI
implementation. For Ubuntu 22.04, the following commands set up the necessary
build tooling, except for MPI.

\begin{bash}
# Install basic utilities
sudo apt install -y gcc g++ make cmake autoconf \
automake libtool pkg-config git

# Install rust toolchain via rustup (non-interactive, stable)
curl --proto '=https' --tlsv1.2 -sSf https://sh.rustup.rs \
| sh -s -- -y --default-toolchain stable
. "$HOME/.cargo/env"
\end{bash}

Now, ORCA can be built via the \code|orca-umbrella| artifact. The artifact
bundles vendored versions of ORCA along with key dependencies and scripts.
Internet access is required for third-party Rust/Python dependencies.

\begin{bash}
# For mini runs, OR_PREFIX can be any path on the filesystem
OR_PREFIX=/tmp/orca-prefix
OR_BUILD_DIR=$OR_PREFIX/build
OR_INSTALL_DIR=$OR_PREFIX/install
mkdir -p $OR_PREFIX $OR_BUILD_DIR $OR_INSTALL_DIR

# Obtain orca-umbrella via git-clone, or via Zenodo artifact.
# $ORCA_CMAKE_ROOT must point to the CMake project root.
# Zenodo zips sometimes add a hash after file/dir names.
git clone https://github.com/pdlfs/orca-umbrella.git \
  $OR_PREFIX/orca-umbrella
ORCA_CMAKE_ROOT=$OR_PREFIX/orca-umbrella
ls $ORCA_CMAKE_ROOT # Must list: CMakeLists.txt, cache, etc.

cd $OR_BUILD_DIR
cmake -DCMAKE_INSTALL_PREFIX=$OR_INSTALL_DIR $ORCA_CMAKE_ROOT
make -j<nproc> # Takes 20-30 mins

# Polars required for Artifact 2, but demo script
# interleaves workflows, so it should be installed upfront.
pip install -r $OR_INSTALL_DIR/scripts/requirements.txt
\end{bash}

Once successfully built, all artifacts will be available in \code|$OR_INSTALL|,
with scripts in \code|$OR_INSTALL/scripts|.

\artcomp

Once built, \code|run_orca_ae.sh| will provide a guided tour of all artifact
functionality. The following commands are provided for reference, but need not
be executed manually.

The script will run 10 timesteps of an AMR code augmented with ORCA under
different modes. A single run involves:
\begin{itemize}
	\item The AMR code, run as 16 ranks of \code|phoebus| with Sedov Blast Wave
	      decks generated from \code|decks/blast_wave_3d.512.pin.in|. The AMR
	      code has been modified to link with ORCA's client module.
	\item ORCA nodes: a single controller and aggregator
	      (\code|controller_main| and \code|aggregator_main|, respectively),
	      running on the same node.
\end{itemize}

A number of tunable parameters are described in \code|defaults.yml|, passed to
ORCA nodes via the env-var \code|ORCA_CFG_YAML|. All parameters can be
overridden using corresponding env-vars, and the scripts override some params
to appropriate values for the AMR code.

The script will demonstrate ORCA analytics and control capabilities. Analytics
are controlled using OrcaFlow (ORCA's SQL-based programmable interface) and
control commands. We explain these using examples.

\begin{enumerate}
	\item \code|bootstrap_cmdseq: set-flow enable-tracers; resume| in
	      \code|defaults.yml| defines two commands supplied by the controller to
	      the simulation at bootstrap. \code{set-flow enable-tracers} activates
	      all telemetry from all available trace streams (currently MPI and
	      Kokkos), and \code|resume| asks the simulation to begin right after.
	      ORCA-enabled simulations initialize in a paused state and start only
	      once they have an initial set of commands.
	\item At any point while the simulation is running, additional commands may
	      be supplied using the \code|cmdrunner_main| binary. The same YAML is
	      used by the binary to locate the controller's IP and port to send
	      commands. E.g. \code|cmdrunner_main "pause"| will pause the simulation.
	\item Flows are SQL-based transforms written as OrcaFlow YAMLs, and sent to
	      the application via command: \\
        \null\hfill\code|set-flow file <path/to/flow.yml>|\hfill\null \\
	      They are committed via TS2PC at runtime, and all pushdowns are applied
	      at the timestep for which the command is committed.
	\item Any active telemetry streams are persisted as timestep-partitioned
	      Parquet or DuckDB at either the controller or the aggregators. The
	      scripts only show the Parquet pathway. Parquet files are flushed at
	      configurable timestep intervals.
	\item Additional probes may be enabled or disabled at any point while the
	      simulation runs.
\end{enumerate}

\artout

As the workflow above demonstrates capabilities, minimal analysis is required.
The following capabilities are demonstrated:
\begin{enumerate}
	\item Basic telemetry collection and persistence as Parquet via aggregators
	      and controller, with minimal in-situ analytics.
	\item An example of SQL-based OrcaFlow analytics, showing pushdown impact on
	      data volume.
	\item Interactive control capabilities while the simulation runs: the script
    demonstrates basic \textsf{PAUSE/RESUME} function, but users are encouraged to
	      try other commands at runtime.
\end{enumerate}

\arteval{2}

\artin

\code|scripts/run_queries.py| is provided to demonstrate basic analyses over
ORCA-generated Parquet. As the format is an industry standard, the goal is basic
validation of ORCA-collected telemetry, and not a comprehensive demonstration
of format capabilities.

\artcomp

The script consists of basic Polars queries emitted against the generated
Parquet. The following telemetry streams will be available:

\begin{example}
parquet/
|-- kokkos_events/
|   `-- ts=<timestep_range>/
|       `-- ranks=<rank_range>.parquet
|-- mpi_collectives/
|   `-- ...
|-- mpi_messages/
|   `-- ...
`-- orca_events/ # ORCA-internal telemetry
\end{example}

\artout

ORCA telemetry is guaranteed to be timestep-partitioned. Each
\code|ts=<beg>_<end>| directory only contains telemetry for that timestep. The
script will demonstrate this briefly.

\arteval{3}

\artin

The artifact consists of a prompt that was provided to Claude Code in an
environment where ORCA was run on a cluster with a known pathology, along with
the conversation transcript. While this artifact is provided for archival
purposes, ORCA's control capabilities are evaluated via manual invocation of
\code|cmdrunner_main| as described above.

\artcomp

\begin{bash}
# Controller is found via ORCA_CFG_YAML,
# or env vars ORCA_CTL_DASHHOST:ORCA_CTL_DASHPORT
cmdrunner_main "<cmd>"
# cmd = pause | resume | status | set-flow <file> etc.
\end{bash}

\artout

When run, commands trigger a timestep-linked two-phase commit to control the
application. Different protocol messages, such as prepare/commit times, current
delta etc. appear on console as the protocol proceeds.

\end{document}
